\section{The free radial flux and the weighted Euler operator}
\label{sec:free-maxwell-radial-flux}

The classification of free vector charges also uses dilation about the
marked insertion. Its generator preserves weighted grade, so its flux
can be used without increasing the spectator capacity. We derive the
normalization from the same stress as in
Section~\ref{sec:free-maxwell-ward-flux}.

Fix a centre $x\in\mathbb R^4$ and the vector field $v_x(y)=y-x$.
For the polarized stress $T_{ab}(f,g)$ used there, differentiation of
$j_a=\sum_b(v_x)_bT_{ab}$ gives the trace of the stress plus its
contracted divergence. The trace is zero in four dimensions. Thus the
current is divergence free wherever both arguments are Maxwell waves.
This is the precise reason dilation is allowed even though its matrix
is not skew-symmetric.

Let $p$ be a linear probe of homogeneous jet order $j\in\mathbb N$,
let $f$ be a Maxwell wave on a neighborhood of the support of a smooth
compact cutoff $\chi$, and assume $\chi=1$ near $x$. The regulator
calculation proving \eqref{eq:maxwell-coordinate-charge} remains valid
for $v_x$: the extra differentiated-drive term is the vanishing stress
trace, including for the regulated closed pulse. Differentiation in
the enclosed pulse centre, with the drive fixed at $v_x$ and the
cutoff fixed, therefore gives
\[
 \mathcal F_\chi(j(v_x;P_{p,x},f))
                         =-p[d\iota_{v_x}f;x].
\]
The differentiation under the collar integral has the same compact,
separated dominator as in \eqref{eq:maxwell-single-charge}.

The Bianchi expansion of the variation is now
\[
 (d\iota_{v_x}f)_{ab}(y)
       =(y-x)\cdot\nabla f_{ab}(y)+2f_{ab}(y).
\]
This is again a Maxwell wave: locally it is the derivative at $s=0$
of $e^{2s}f(x+e^s(y-x))$, and the Maxwell and Bianchi differential
equations are preserved by that rescaling. Its homogeneous Taylor
coefficient of order $j$ at $x$ is $(j+2)$ times that of $f$.
Indeed the degree-$j$ coefficient of the rescaled wave is
$e^{(j+2)s}$ times the original coefficient; differentiating gives
the factor, including for the normalized Taylor coefficient $D^jf/j!$.
Consequently
\begin{equation}\label{eq:maxwell-radial-single-charge}
 \mathcal F_\chi(j(v_x;P_{p,x},f))=-(j+2)p[f;x].
\end{equation}

Write $X_{j,\alpha}$ for the polynomial jet coordinates and define the
weighted Euler derivation
\[
 \mathscr E=\sum_{j,\alpha}(j+2)X_{j,\alpha}
                         \frac{\partial}{\partial X_{j,\alpha}}.
\]
The sum acts finitely on every polynomial. On a monomial it multiplies
by its weighted degree, and hence preserves every finite weighted
filtration. It preserves the gauge-invariant polynomial space since
colour rotations mix only coordinates of the same jet order.

\begin{proposition}[Radial Ward identity for all polynomial words]
\label{prop:maxwell-radial-ward}
Let $x_0,\ldots,x_n$ be distinct points, and let $P_i$ be arbitrary
Wick polynomials of Maxwell jets. Choose $\chi\in C_c^\infty$ equal
to one near $x_0$, with support avoiding $x_i$ for $i>0$. For the
same Wick stress $\mathsf T_{ab}$ as before,
\begin{equation}\label{eq:maxwell-radial-polynomial-ward}
 \mathcal F_\chi\!\left(
   y\longmapsto\Bigl[\sum_b(y-x_0)_b
       \mathcal M_{(y,x_0,\ldots,x_n)}
          (\mathsf T_{ab},P_0,\ldots,P_n)\Bigr]_a\right)
 =-\mathcal M_{(x_0,\ldots,x_n)}
                (\mathscr E P_0,P_1,\ldots,P_n).
\end{equation}
All integrals on the left are absolutely convergent.
\end{proposition}

\begin{proof}
The cutoff derivative has compact support separated from every insertion;
the integrand is smooth there. This proves absolute integrability and
allows the finite Wick sums to pass through the integral. Extract the
two factors in the quadratic stress, as in the proof of
\eqref{eq:maxwell-polynomial-ward}. If neither contracts with $P_0$,
the associated two pulses are smooth throughout the cutoff support,
and the divergence-free current has zero flux. If both contract with
$P_0$, collar independence and pulse homogeneity give
$Z=t^{1-5-j-l}Z$ for jet orders $j,l$ of those two legs: the radial
drive is homogeneous of degree one. Taking $t=2$ forces $Z=0$.
This argument uses only the compact collar, not a product of singular
pulses through their common centre.

If exactly one stress leg contracts with $P_0$, apply
\eqref{eq:maxwell-radial-single-charge} to the other, smooth pulse.
Each occurrence of an order-$j$ coordinate in $P_0$ acquires the factor
$-(j+2)$. Summing occurrences is exactly the action of $-\mathscr E$.
The two choices of stress leg cancel the factor $1/2$ in the polarized
Wick stress, with the same normalization as the preceding polynomial
Ward identity. The remaining Wick pairings are those of the word on
the right of \eqref{eq:maxwell-radial-polynomial-ward}. Linearity in
all polynomials completes the proof.
\end{proof}
